% Part: sets-functions-relations
% Chapter: infinite
% Section: dedekind-induction

\documentclass[../../../include/open-logic-section]{subfiles}

\begin{document}

\olfileid{sfr}{infinite}{induction}
\olsection[Arithmetical Induction]{డెడెకిండ్ బీజగణితాలు, అంకగణిత ఆగమనం}

ఇప్పుడు కీలకంగా, ఒక డెడెకిండ్ బీజగణితం---నిజానికి \emph{ఏ}
డెడెకిండ్ బీజగణితమైనా---సహజ సంఖ్యలకు ప్రతిరూపంగా పనిచేస్తుంది.
కింది సులభ పరిణామమే దీనికి కారణం:

\begin{thm}[అంకగణిత ఆగమనం]\ollabel{thm:dedinfiniteinduction}
$N, s, o$ కలిసి ఒక డెడెకిండ్ బీజగణితాన్ని ఏర్పరుస్తాయనుకుందాం.
అప్పుడు ఏ సమితి $X$కైనా:
	\begin{center}
		if $o \in X$ and $(\forall {n} \in N \cap X){s}({n}) \in X$, {then} $N \subseteq X$.
	\end{center}
\end{thm}

\begin{proof}
డెడెకిండ్ బీజగణిత నిర్వచనం ప్రకారం
$N = \closureofunder{s}{o}$. ఇప్పుడు ${o} \in X$, అలాగే
$(\forall {n} \in N)(n \in X \lif {s}({n}) \in X)$ అయితే,
$N \cap X$లో $o$ ఉంటుంది, అది $s$-సంవృతం. కాబట్టి
$N = \closureofunder{s}{o} \subseteq N \cap X \subseteq X$.
\sourcecorrection{OLTEINF-003}{ఆంగ్ల మూలం ఏ సమితి $X$నైనా అనుమతించి,
స్వ-ప్రమేయపు పరిధిలోని ఉపసమితులకే వర్తించే సంవృత కనిష్ఠతను నేరుగా
$X$కు వర్తింపజేసింది. సిద్ధాంతంలోని షరతే $N\cap X$లో $o$ ఉండి అది
$s$-సంవృతమని ఇస్తుంది; ఇక్కడ కనిష్ఠతను ఆ సమితికి వర్తింపజేశాం.}
\end{proof}

ఆగమనం సహజ సంఖ్యల విశిష్ట ధర్మం కనుక, విషయం ఇది. ఏ డెడెకిండ్ అనంత
సమితి ఇచ్చినా, ఒక డెడెకిండ్ బీజగణితాన్ని ఏర్పరచి, దానిని సహజ
సంఖ్యలకు మన ప్రతిరూపంగా ఉపయోగించవచ్చు.

\olref{thm:dedinfiniteinduction}, ఆగమనాన్ని \emph{సమితి-సిద్ధాంత}
పదాల్లో సూత్రీకరిస్తుందని ఒప్పుకోవాలి. కానీ ఈ సూత్రాన్ని మరింత
పరిచితంగా ఉండే పదాల్లో సులభంగా పెట్టవచ్చు:

\begin{cor}\ollabel{natinductionschema}
$N, s, o$ కలిసి ఒక డెడెకిండ్ బీజగణితాన్ని ఏర్పరుస్తాయనుకుందాం. అప్పుడు
పరామితులు ఉండగల ఏ సూత్రం $\phi(x)$కైనా:
\begin{center}
	if $\phi(o)$ and $(\forall {n} \in N)(\phi(n)\lif
	\phi({s}({n})))$, {then} $(\forall n \in N)\phi(n)$
\end{center}
\end{cor}

\begin{proof}
$X = \Setabs{n \in N}{\phi(n)}$గా తీసుకుని, ఇప్పుడు
\olref{thm:dedinfiniteinduction}ను ఉపయోగించండి.
\end{proof}

ఈ ఫలితంలో ఒక సూత్రానికి ``పరామితులు ఉండటం'' గురించి మాట్లాడాం.
స్థూలంగా దాని అర్థం ఏమిటంటే, ఏ వస్తువులు $c_1, \ldots, c_k$కైనా
$\phi(x, c_1, \ldots, c_k)$తో పనిచేయవచ్చు. మరింత కచ్చితంగా,
``పరామితులు'' అని ప్రస్తావించకుండా ఫలితాన్ని ఇలా ప్రకటించవచ్చు.
స్వేచ్ఛా చరరాశులన్నీ చూపబడిన ఏ సూత్రం
$\phi(x, v_1, \ldots, v_k)$కైనా:
	\begin{align*}
			\forall v_1 \ldots \forall v_k((&\phi(o, v_1,\ldots, v_k) \land {}\\
			&	(\forall x \in N)(\phi(x,v_1, \ldots, v_k) \lif \phi(s(x), v_1,\ldots, v_k))) \lif {}\\
			&\hspace{3em} (\forall x \in N)\phi(x, v_1,\ldots, v_k))
	\end{align*}
``పరామితులు ఉండటం'' అని మాట్లాడటం చదవడాన్ని చాలా సులభం చేస్తుందని
స్పష్టం. (\olref[sth][][]{part}లో ఈ పద్ధతిని తరచుగా ఉపయోగిస్తాం.)

మళ్లీ డెడెకిండ్ బీజగణితాలకు వద్దాం: ఏ డెడెకిండ్ బీజగణితం ఇచ్చినా,
సంకలనం, గుణకారం, ఘాతాంకం అనే సాధారణ అంకగణిత ప్రమేయాలను కూడా
నిర్వచించవచ్చు. అయితే ఇది సులభమైన విషయం కాదు; ఇందులో
\emph{పునరావృత్త నిర్వచనం} అనే పద్ధతి ఉంటుంది. ఆ పద్ధతిని చాలా
తరువాత, మరింత సాధారణ సందర్భంలో పరిచయం చేసి సమర్థిస్తాం.
(ఆసక్తిగలవారు \olref[sth][ord-arithmetic][]{chap} తరువాత దీనిని
మళ్లీ పరిశీలించవచ్చు, లేదా \citealt[pp.~95--8]{Potter2004} వంటి
ప్రత్యామ్నాయ వివరణను చదవవచ్చు.) అయితే $N, s, o$ కలిసి ఒక డెడెకిండ్
బీజగణితాన్ని ఏర్పరచినప్పుడు, చివరికి కింది వాటిని నిర్దేశించగలం:
\begin{align*}
	{a} + {o} &= {a} & & & {a} \times {o} &= {o} & & & {a}^{o} &= s(o)\\
	{a} + {s}({b}) &= {s}({a}+{b}) &&& {a} \times {s}({b}) &= ({a}\times {b}) + {a}  & & & {a}^{{s}({b})} &= {a}^{b} \times {a}
\end{align*}
ఇవి మనం ఆశించే విధంగానే ప్రవర్తిస్తాయని చూపగలం.

\end{document}
